\documentclass[11pt]{article}
\usepackage[a4paper,margin=1in]{geometry}
\usepackage{fontspec}
\usepackage[boldfont=false]{xeCJK}
\setCJKmainfont{FandolSong-Regular.otf}
\usepackage{amsmath}
\usepackage{amssymb}
\usepackage{array}
\usepackage{booktabs}
\usepackage{graphicx}
\usepackage{adjustbox}
\usepackage{enumitem}
\setlist[itemize]{topsep=2pt,partopsep=0pt,parsep=0pt,itemsep=2pt,leftmargin=1.4em}
\setlist[enumerate]{topsep=2pt,partopsep=0pt,parsep=0pt,itemsep=2pt,leftmargin=1.6em}
\usepackage{parskip}
\usepackage{fvextra}
\fvset{breaklines=true,breakanywhere=true,fontsize=\small}
\setlength{\parindent}{0pt}

\begin{document}

\begin{center}
  {\LARGE\bfseries Further Programming}\\[0.35em]
  {\large A Level Computer Science}
\end{center}
\vspace{0.6em}

\section*{Programming paradigms}

A \textbf{programming paradigm} 编程范式 is a style of programming — a way of structuring programs, with its own ideas and language features. Four are in this syllabus.

\subsection*{Low-level programming}

Programming \textbf{close to the hardware} in \textbf{machine code} 机器码 or \textbf{assembly language} 汇编语言, where each instruction maps to what the CPU runs. It gives direct access to \textbf{registers} 寄存器 and \textbf{memory addresses} 内存地址, and is very fast and compact, but is architecture-specific, tedious, and hard to maintain. This is \textbf{low-level} 低级 programming. Used for device drivers, firmware and bootloaders.

\subsection*{Imperative (procedural) programming}

In \textbf{imperative programming} 命令式编程 the programmer writes a \textbf{sequence of commands} that change the program's state — assignments, conditionals, loops, function calls. \textbf{Variables} 变量 hold state; statements change it; code is organised into procedures and functions. This is the style of Topics 9 and 11 (Python, C). Strong when the algorithm has clear sequential steps.

\subsection*{Object-oriented programming (OOP)}

In \textbf{object-oriented programming} 面向对象编程 programs are built from \textbf{objects} 对象 — units combining \textbf{data} (\textbf{attributes} 属性) and \textbf{operations} (\textbf{methods} 方法). Objects are \textbf{instances} 实例 of \textbf{classes} 类. The four pillars:

\begin{itemize}
\item \textbf{encapsulation} 封装 — an object's data is hidden behind its methods; outside code uses the public methods only, not the data directly. This protects the object and lets its internals change without breaking callers.

\item \textbf{inheritance} 继承 — a \textbf{subclass} 子类 specialises a \textbf{superclass} 父类, inheriting its attributes and methods and adding or \textbf{overriding} 重写 them. Models "is-a" ("a Manager is an Employee").

\item \textbf{polymorphism} 多态 — different objects respond to the \textbf{same method call} differently; the caller need not know the exact type. Every \texttt{Shape} has \texttt{Area()}, and a \texttt{Circle} and a \texttt{Rectangle} each implement it their own way.

\item \textbf{abstraction} 抽象 — show a simple interface and hide the implementation.

\end{itemize}

Other terms: a \textbf{constructor} 构造函数 is a special method run when an object is created, to set up its attributes. Used for large systems, GUIs, simulations and games.

\par\medskip
\noindent\begin{minipage}{\linewidth}
\centering
\adjustbox{max width=\linewidth, max totalheight=0.42\textheight}{%
  \includegraphics{/Users/joschka/Desktop/Programming/eXercise/output/eXam/handouts/a_level_computer_science/assets/20-p517-class-diagram.pdf}}\\[0.3em]
\emph{A class diagram for a Shape: private attributes and public methods}
\end{minipage}\par\medskip

\par\medskip
\noindent\begin{minipage}{\linewidth}
\centering
\adjustbox{max width=\linewidth, max totalheight=0.42\textheight}{%
  \includegraphics{/Users/joschka/Desktop/Programming/eXercise/output/eXam/handouts/a_level_computer_science/assets/20-p525-inheritance.pdf}}\\[0.3em]
\emph{Inheritance: partTime and fullTime are subclasses of employee}
\end{minipage}\par\medskip

\subsection*{Declarative programming}

In \textbf{declarative programming} 声明式编程 you say \textbf{what} to compute, not \textbf{how} — the runtime works out the steps. Two kinds:

\begin{itemize}
\item \textbf{functional programming} 函数式编程 — built from \textbf{pure functions} 纯函数 (no \textbf{side effects} 副作用; same input always gives the same output) composed together. Examples: Haskell, Lisp.

\item \textbf{logic programming} 逻辑编程 — state facts and rules; the engine answers queries by inference. Example: Prolog.

\end{itemize}

A familiar declarative example is \textbf{SQL} 结构化查询语言: \texttt{SELECT * FROM Customer WHERE Country = 'UK'} says what you want, not how to walk the records.

\subsection*{Comparing paradigms}

\begin{center}
\adjustbox{max width=\linewidth}{%
\begin{tabular}{lll}
\toprule
\textbf{Paradigm} & \textbf{Strength} & \textbf{Typical languages} \\
\midrule
Low-level & maximum control, speed & assembly \\
Imperative & direct, intuitive & C, Python \\
Object-oriented & modular, models entities & Java, C\#, Python \\
Functional & clear, no side effects & Haskell, F\# \\
Logic & inference, rules & Prolog \\
Database & data queries & SQL \\
\bottomrule
\end{tabular}}
\end{center}

Modern languages often \textbf{mix paradigms} — Python supports all of procedural, OOP and functional. The right one depends on the problem.

\section*{File processing}

This extends the \textbf{file} 文件 handling from Topic 10. Pseudocode operations: \texttt{OPENFILE name FOR READ | WRITE | APPEND} (READ opens an existing file, WRITE creates/overwrites, APPEND adds to the end); \texttt{READFILE name, line}; \texttt{WRITEFILE name, value}; \texttt{CLOSEFILE name}; and \texttt{EOF(name)} which is TRUE at the end.

Read a whole file:

\begin{Verbatim}
OPENFILE "names.txt" FOR READ
WHILE NOT EOF("names.txt") DO
    READFILE "names.txt", thisName
    OUTPUT thisName
ENDWHILE
CLOSEFILE "names.txt"
\end{Verbatim}

Search a file (stop when found):

\begin{Verbatim}
found ← FALSE
OPENFILE "people.txt" FOR READ
WHILE NOT EOF("people.txt") AND NOT found DO
    READFILE "people.txt", line
    IF line = target THEN found ← TRUE
ENDWHILE
CLOSEFILE "people.txt"
\end{Verbatim}

\subsection*{Updating a file in place}

Most languages can't edit a text file in place. Instead: open the original for READ and a temporary file for WRITE; for each line, write the new version if it should change, else the original; close both; then replace the original with the temp file. The same pattern handles deleting lines (skip them) and inserting lines.

\subsection*{Pitfalls}

Forgetting to close a file (data may be lost); opening for WRITE when you meant APPEND (overwrites everything); reading past \texttt{EOF}; hard-coded paths (make them constants for portability).

\section*{Exception handling}

An \textbf{exception} 异常 is an error or unexpected condition during execution — divide by zero, file not found, network failure, an \textbf{array} 数组 index out of range. \textbf{Exception handling} 异常处理 lets a program \textbf{detect} it and respond gracefully instead of crashing.

It matters because real programs face errors that \textbf{cannot be prevented up front} (files moved, networks down, bad input); without it, every operation needs its own \texttt{IF} check; and it \textbf{separates} the normal flow from the error handling, so the main path reads cleanly.

\subsection*{Pattern}

\begin{Verbatim}
TRY
    OPENFILE "data.txt" FOR READ
    READFILE "data.txt", line
    OUTPUT line
    CLOSEFILE "data.txt"
EXCEPT FileNotFound
    OUTPUT "Sorry, the file does not exist."
EXCEPT ReadError
    OUTPUT "Sorry, error reading the file."
ENDTRY
\end{Verbatim}

The \texttt{TRY} block holds the code that might fail; the first matching \texttt{EXCEPT} block runs. Real languages also have a catch-all \texttt{EXCEPT} and a \texttt{FINALLY} block that runs whether or not an exception happened — useful for cleanup (closing files).

\subsection*{Raising an exception}

A subroutine that detects an error can \textbf{raise} 抛出 an exception so the caller handles it:

\begin{Verbatim}
PROCEDURE Divide(a : INTEGER, b : INTEGER) RETURNS INTEGER
    IF b = 0 THEN
        RAISE DivideByZero
    ENDIF
    RETURN a DIV b
ENDPROCEDURE
\end{Verbatim}

\subsection*{Where to handle exceptions}

Handle them \textbf{close to the error} if the response is simple (a message, a retry), or \textbf{higher up the call stack} 调用栈 if only the outer code knows what to do (a top-level GUI loop logs the error and shows a friendly dialog). \textbf{Don't swallow exceptions silently} — at least log them, or debugging becomes impossible.

Common exceptions: \texttt{FileNotFound}, \texttt{IOError}, \texttt{DivisionByZero}, \texttt{IndexOutOfRange}, \texttt{InvalidArgument}, \texttt{NullReference}, \texttt{OutOfMemory}. Wrapping each failing operation in a \texttt{TRY} with the right \texttt{EXCEPT} handlers gives a program that \textbf{degrades gracefully} instead of crashing.


\end{document}
